% English translation of chapter2.tex lines 9--32.
% Source: Methods of Algebra, Volume 2, commit
% 9a5803ff77dd3257484cb177f851a73770a59dd3 (CC BY 4.0).
% stable-unit-id: o014.aljabr2.chapter2.overview

% segment-id: o014.li2.chapter2.overview.h001
\chapter{Abelian Categories}\label{sec:Abel-cat}
\hypertarget{o014.aljabr2.chapter2.overview}{ }

% segment-id: o014.li2.chapter2.overview.p001
An abelian category is an additive category with kernels and cokernels in
which every morphism is strict. Abelian categories generalize the category
$R\dcate{Mod}$ of left modules or $\cated{Mod}R$ of right modules over a ring
$R$, and support the study of such notions as injective and projective
objects, complexes, homology and cohomology, and exact sequences.

% segment-id: o014.li2.chapter2.overview.p002
Module categories are the prototype of abelian categories, but do not exhaust
their substance. There are two reasons for this.
% segment-id: o014.li2.chapter2.overview.l001
\begin{enumerate}
% segment-id: o014.li2.chapter2.overview.i001
	\item In a general abelian category, objects have no elements and
	morphisms are not maps. Diagram-chasing methods from module categories are
	therefore no longer available; the most basic exact sequences must instead
	be derived by operations on objects and functors. Prominent examples are
	the snake lemma and the five lemma discussed in
	\S\ref{sec:diagram-lemmas}, whose proofs use a substitute for diagram
	chasing (Lemma \ref{prop:Abel-cat-exact-aux}). Even so, the arguments are
	considerably more involved than their module-category counterparts.

% segment-id: o014.li2.chapter2.overview.i002
	\item Even when one starts with a module category or one of its abelian
	subcategories, abstract abelian categories arise when studying functor
	categories, localizations, or Serre quotients. In geometry and topology,
	sheaf theory naturally leads to such constructions.
\end{enumerate}

% segment-id: o014.li2.chapter2.overview.p003
The Freyd--Mitchell embedding theorem \ref{prop:FM} states that every small
abelian category admits a fully faithful embedding into some
$\cated{Mod}R$, and that the embedding functor preserves exact sequences.
Once this result is accepted, many assertions about exact sequences reduce
readily to the case of $\cated{Mod}R$. This reduction keeps diagram chasing
available, but its chief benefit may be psychological reassurance: on the one
hand it sidesteps the essence of abelian categories, and on the other hand any
proof of the embedding theorem must still establish several core facts about
abelian categories from scratch. This chapter does not depend on the
embedding theorem, nor logically on module theory, but readers are expected
to have a basic knowledge of homological algebra in module categories.
Diagram-chasing skills remain useful and may even be indispensable in the
natural order of study; readers are encouraged to consult \cite{Li1} or the
relevant parts of another textbook. Appendix \S\ref{sec:FM} derives the
embedding theorem as an application of the Ind construction.

% segment-id: o014.li2.chapter2.overview.p004
A cocomplete abelian category with a generator in which filtered small
colimits are exact is called a Grothendieck category. Such categories are the
subject of the final section, \S\ref{sec:Grothendieck-cat}; these conditions
are directly related to assumptions on set size. Grothendieck categories
occur frequently in geometry and enjoy better properties, such as
completeness (Corollary \ref{prop:Grothendieck-cat-complete}) and the existence
of canonical injective resolutions (Theorem
\ref{prop:Grothendieck-injective}). They are a closer analogue of module
categories than general abelian categories; the precise meaning of this
statement is explained by the Gabriel--Popescu theorem \ref{prop:GP} in the
appendix.

% segment-id: o014.li2.chapter2.overview.p005
The basic isomorphism theorems of module theory remain valid for general
abelian categories. The same is true of such fundamental notions as
subobjects, quotient objects, direct-sum decompositions, indecomposable
objects, simple objects, and composition series, all central concerns of
algebra. They will be treated in \S\S\ref{sec:direct-sum}---
\ref{sec:semisimple}. For example, the Krull--Remak--Schmidt theorem on
direct-sum decompositions has an abstract generalization (Theorem
\ref{prop:Krull-Schmidt-gen}). The language of lattice theory is useful here;
see \S\ref{sec:lattices}. It should be stressed that statements involving
infinite direct sums often hold only in Grothendieck categories.

% segment-id: o014.li2.chapter2.overview.p006
Section \ref{sec:Serre-subcat} begins with the definition and characterization
of abelian subcategories, and then introduces quotients of an abelian category
by a certain kind of abelian subcategory, called a Serre subcategory; the
result is a Serre quotient. These quotients are characterized by a universal
property and can be constructed by the localization introduced in
\S\ref{sec:cat-localization}. The section also introduces the group
$\mathrm{K}_0$ of an abelian category and its basic properties, including the
Euler--Poincar\'e principle (Theorem \ref{prop:EP}) and the exact sequence for
$\mathrm{K}_0$ of a Serre quotient (Proposition
\ref{prop:K0-exact-seq}). These techniques are common in applications. The
construction of $\mathrm{K}_0$ applies to a broader class of structures than
abelian categories, called \emph{exact categories}; see \cite{Bu10} for the
theory. Moreover, $\mathrm{K}_0$ is only the beginning of a sequence of
abelian groups $\mathrm{K}_n$ for $n\in\Z_{\geq0}$. A satisfactory
interpretation of the higher K-groups requires homotopy-theoretic ideas,
especially the notion of a ``spectrum'', and lies beyond the scope of this
book.
\index{exact category}

% segment-id: o014.li2.chapter2.overview.n001
\begin{petunjukbacaan}
	As noted above, readers should have some familiarity with concrete
	operations in module categories. For the basic theory of complexes and
	their cohomology, the material in
	\S\S\ref{sec:Abel-cat-def}---\ref{sec:inj-proj} is essential. Apart from
	the treatment of abelian subcategories, the material in
	\S\ref{sec:Serre-subcat} is used less often later, but it remains standard
	knowledge and is related to several definitions in the discussion of
	derived categories. Section \ref{sec:Grothendieck-cat} on Grothendieck
	categories is supplementary; it requires relatively involved arguments
	and material from the appendices, and may be skipped on a first reading.

	A complex written as data $(X^n,d^n)_n$, where
	$d^n:X^n\to X^{n+1}$, is also called a cochain complex. If instead one
	uses subscripts, writing $(X_n,d_n)_n$ with
	$d_n:X_n\to X_{n-1}$, it is called a chain complex. Except in
	Chapter~\sourcecrossref{sec:simplicial}{8}, this book primarily uses
	cochain notation.
\end{petunjukbacaan}
